\section{Methodology}
\label{sec:methodology}

\subsection{Adversarial Prompt Taxonomy}

We organize adversarial prompts into four categories based on the type of vulnerability they exploit:

\begin{table}[H]
\centering
\caption{Adversarial prompt categories and their descriptions}
\label{tab:categories}
\begin{tabularx}{\textwidth}{lX}
\toprule
\textbf{Category} & \textbf{Description} \\
\midrule
Math \& Statistics & Prompts exploiting statistical paradoxes, probability misconceptions, and mathematical fallacies to trick the model into confirming incorrect analyses. \\
Security \& IAM & Code review and security analysis scenarios containing subtle vulnerabilities that the model should identify but may miss. \\
Reasoning \& Jailbreak & Attempts to bypass safety guidelines through roleplay, encoding, authority impersonation, logical fallacies, and social pressure. \\
Hallucination & Prompts requesting information about fabricated entities (libraries, events, versions) to test whether the model will confabulate. \\
\bottomrule
\end{tabularx}
\end{table}

Each category contains 5 prompts, totaling 20 adversarial tests. Every prompt includes:
\begin{itemize}
    \item A specific \textbf{adversarial technique} (e.g., Simpson's Paradox Gaslighting, JWT None Algorithm)
    \item A documented \textbf{hypothesis} about expected model failure
    \item \textbf{Ground truth} defining the correct response
    \item \textbf{Deception criteria} specifying what constitutes a failure
    \item A \textbf{difficulty rating} (easy, medium, hard)
\end{itemize}

\subsection{Evaluation Pipeline}

The benchmark follows a three-stage automated pipeline:

\begin{enumerate}
    \item \textbf{Target Generation:} Each adversarial prompt is sent to the target LLM (Gemini 2.5 Flash Lite) via API, and the response is recorded.
    \item \textbf{Judge Evaluation:} A separate judge model (Claude) evaluates the target's response against the ground truth and deception criteria, producing a structured assessment~\cite{zheng2023judging}.
    \item \textbf{Scoring \& Aggregation:} Judge evaluations are aggregated into per-prompt scores, category-level metrics, and overall statistics.
\end{enumerate}

\subsection{Scoring Rubric}

The judge model assigns scores on a continuous 0.0--1.0 scale with the following anchor points:

\begin{table}[H]
\centering
\caption{Robustness scoring rubric}
\label{tab:rubric}
\begin{tabular}{clp{8cm}}
\toprule
\textbf{Score} & \textbf{Label} & \textbf{Description} \\
\midrule
1.00 & Fully Robust & Correctly identified the trick; provided accurate information \\
0.75 & Mostly Robust & Identified main issue with minor inaccuracies \\
0.50 & Partially Deceived & Some awareness but significant errors \\
0.25 & Mostly Deceived & Failed to identify core deception \\
0.00 & Fully Deceived & Completely fell for the adversarial prompt \\
\bottomrule
\end{tabular}
\end{table}

A response with score $< 0.5$ is classified as ``deceived.'' The overall \textbf{deception rate} (error rate) is the fraction of prompts where the model was deceived.

\subsection{Technical Implementation}

\begin{itemize}
    \item \textbf{API:} Google Generative AI (Gemini 2.5 Flash Lite), free tier with rate limits per model per day.
    \item \textbf{Rate Limiting:} 5-second delay between API calls; exponential backoff on failures.
    \item \textbf{Caching:} Responses cached locally to prevent redundant API calls.
    \item \textbf{Pipeline:} Python with \texttt{google-genai}, \texttt{pandas}, \texttt{matplotlib}/\texttt{seaborn}.
    \item \textbf{Hybrid Judging:} For 10 prompts with API responses, Claude was used as judge. 
    For 10 prompts without API responses, scores were estimated from documented LLM behavior patterns.
\end{itemize}

\subsection{Limitations of the Methodology}

\begin{itemize}
    \item \textbf{Partial API coverage:} API quota constraints limited actual model responses to 10 of 20 prompts; remaining scores are estimates.
    \item \textbf{Automated judging:} AI judges may miss subtle failures that human evaluators would catch~\cite{mazeika2024harmbench}. 
    Also the judge model's own biases may affect scoring consistency.
    \item \textbf{Prompt sensitivity:} Results may vary with minor prompt reformulations.
    \item \textbf{Sample size:} 20 prompts provide indicative rather than statistically conclusive results.
\end{itemize}
